% ============================================================
% SECTION 3 — Spiritual Resistance (Chapters 3 & 5)
% ============================================================
\section{Spiritual Resistance}

\begin{frame}{Zayn al-'Abidin: Worship as a Journey of the Soul}
\begin{itemize}[itemsep=5pt]
  \item During his worship, his child broke a bone -- the household treated the child without interrupting him; he noticed only the next morning
  \item ``Herald of affection'' -- deliberately sought out the socially invisible, even inviting lepers into his home
  \item On Hajj, chose anonymity specifically \emph{to be able to serve} the caravan, not be served by it
\end{itemize}
\bigquote{[I chose unfamiliarity] to gain the opportunity to serve Muslims and friends.}{Zayn al-'Abidin}
\delivernote{Don't resolve the tension in the broken-bone story for the room -- let them feel both readings: genuine devotion, and a real question about what a father owes a child.}
\end{frame}

\begin{frame}{Mourning as Purposeful Remembrance}
\begin{itemize}[itemsep=5pt]
  \item Decades of visible grief for al-Husayn -- not unresolved trauma, but deliberate preservation of \emph{why} he rose and \emph{who} bore responsibility
  \item Compared to Ya'qub's grief for Yusuf, in depth
  \item Service to Islam takes different valid forms -- prayer and supplication are Zayn al-'Abidin's own legitimate legacy, not a lesser substitute for the sword
\end{itemize}
\end{frame}

\begin{frame}{Musa al-Kazim: Four Years, Dignity Under Pressure}
\begin{itemize}[itemsep=5pt]
  \item Repeatedly offered freedom for confession or submission to Harun -- repeatedly refused
  \item Jailers 'Isa ibn Ja'far, Fadl ibn Rabi', Fadl ibn Yahya al-Barmaki all grow sympathetic through direct contact -- Harun keeps transferring him and punishing sympathetic jailers
  \item Near the end: offered release for a simple, private confession -- refused, sensing death was near
\end{itemize}
\bigquote{Time can never again bring a person like 'Ali.}{Mu'awiyah, reportedly, decades earlier -- the same pattern of grudging private respect recurs with Harun}
\end{frame}

\begin{frame}{Influence Without a Platform}
\begin{itemize}[itemsep=5pt]
  \item \textbf{Bishr Hafi:} one indirect question (``does this house belong to a free man or a slave?'') --- reformed without a single word of rebuke
  \item \textbf{Safwan Jammal:} condemned for quietly hoping Harun would live long enough to pay a camel-rental debt -- self-interest entangled with an oppressor's survival
  \item \textbf{'Ali ibn Yaqtin:} encouraged to \emph{remain} inside Harun's government, concealing his Shi'ism, because the position let him protect other believers
\end{itemize}
\examplenote{The dividing line}{Cooperation with unjust power is judged by whose interest it serves -- self-interest condemns it (Safwan); service to the community can justify it (Ibn Yaqtin).}
\end{frame}

\begin{frame}{Three-Part Explanation for the Martyrdoms}
\begin{enumerate}[label=\arabic*.,itemsep=6pt]
  \item Spiritual authority itself threatened ruling caliphs -- no army required
  \item The Imams opposed oppression while practicing taqiyyah, leaving no usable evidence against them
  \item An exceptional spirit of resistance that could not be induced to surrender
\end{enumerate}
\delivernote{Ask the room: how are an Imam practicing taqiyyah (al-Sadiq, next section) and an Imam refusing all compromise (Musa al-Kazim, here) not actually contradicting each other? (Both are the same underlying commitment -- refusing to legitimize illegitimate power -- applied to different circumstances.)}
\end{frame}

\qaframe{Spiritual Resistance}{%
  \qa{What single concept links a prisoner with no political power (Musa al-Kazim) to a crown prince inside the government (al-Rida, next section)?}
     {Spiritual authority is a distinct, ruler-uncontrollable currency of legitimacy -- it needs no army, office, or even physical freedom to threaten a caliph's hold on power, because it draws people's ultimate loyalty toward something the ruler doesn't control.}
  \qa{Bishr Hafi and Safwan Jammal both show intervention through a single pointed question rather than a lecture. What does this teach about how the text wants ``religious influence'' understood?}
     {Influence operates person by person, through minimal, well-aimed interventions that require the listener to complete the reasoning themselves -- contrasted explicitly with organized propaganda or open confrontation.}
  \qa{Compare Musa al-Kazim's refusal of Harun's confession offer with al-Sadiq's refusal to even read Abu Salmah's letter (next section). What's actually different?}
     {Al-Sadiq refuses an offer of power before any deal is struck -- he loses an opportunity he never firmly had. Musa al-Kazim refuses a demand for submission from a power that already holds him captive -- he loses his life. Same principle, much higher stakes in the second case.}
}
